% !TEX root = ../FundationsDataScience.tex

%%%%%%%%%%%%%%%%%%%%%%%%%%%%%%%%%%%%%%%%%%%%%%%%%%%%%%%%%%%%%%%%%%%%%%%%
%%%%%%%%%%%%%%%%%%%%%%%%%%%%%%%%%%%%%%%%%%%%%%%%%%%%%%%%%%%%%%%%%%%%%%%%
%%%%%%%%%%%%%%%%%%%%%%%%%%%%%%%%%%%%%%%%%%%%%%%%%%%%%%%%%%%%%%%%%%%%%%%%
\section{Stochastic Optimization}
\label{sec-stochastic-optim}

We study stochastic gradient methods for large datasets of size $n$ and for objectives defined as expectations.

% We set the classes indexes to be $\{-1,+1\}$, and remove empty features, normalize $X$. $n$ is the number of samples, $p$ is the dimensionality of the features,

%%%%%%%%%%%%%%%%%%%%%%%%%%%%%%%%%%%%%%%%%%%%%%%%%%%%%%%%
\subsection{Minimizing Sums and Expectations}

Many learning objectives are finite averages:
\eql{\label{eq-min-sums}
	\umin{x \in \RR^p} f(x) \eqdef \frac{1}{n} \sum_{i=1}^n f_i(x)
}
or expectations:
\eql{\label{eq-min-int}
	\umin{x \in \RR^p}  f(x) \eqdef \EE_{\zp \sim \pi}( f(x,\zp) ) = \int_{\Zz} f(x,z) \d\pi(z).
}
A finite average is an expectation under the uniform measure $\pi=\frac1n\sum_{i=1}^n\de_i$, with $f(x,i)=f_i(x)$. Conversely, drawing $(z_i)_i$ i.i.d. from $\pi$ and setting $f_i(x)=f(x,z_i)$ gives an empirical approximation of~\eqref{eq-min-int}. Under integrability, the empirical objective converges pointwise by the law of large numbers. Convergence of minima or minimizers requires additional assumptions.

For example, empirical risk minimization with linear predictors uses
\eql{\label{eq-stochastic-erm}
	f_i(x) = \ell(\dotp{a_i}{x},y_i)
	\qandq
	f(x,z) = \ell(\dotp{a}{x},y)
}
for $z=(a,y) \in \Zz = (\Aa=\RR^p) \times \Yy$ (typically $\Yy=\RR$ or $\Yy=\{-1,+1\}$ for regression and classification), where $\ell$ is some loss function. 
%
We illustrate the methods with binary logistic classification, using
\eql{\label{eq-stoch-logistic}
	\loss(s,y) \eqdef \log( 1+\exp(-sy) ).
}
The same computational approach extends to general parametric models, including deep neural networks.

The stochastic methods below apply both to finite sums~\eqref{eq-min-sums} and to expectations~\eqref{eq-min-int}, under the corresponding unbiasedness and moment assumptions. We present them for finite sums. When $n$ is large, the aim is to make each iteration inexpensive without evaluating every summand.

When the summands $f_i(x)$ are similar, a single gradient can approximate the average well: $\nabla f_i \approx \nabla f$. Evaluating $\nabla f_i$ is typically $n$ times cheaper than forming the full gradient.
%
Stochastic optimization is nevertheless not always faster than batch gradient descent.
%
For moderate $n$, a deterministic method may use the available computation more effectively.
%
When $n$ is too large for a full pass to be practical, stochastic methods divide the work into small updates and begin improving the model immediately. Mini-batches also permit parallel computation.


%%%%%%%%%%%%%%%%%%%%%%%%%%%%%%%%%%%%%%%%%%%%%%%%%%%%%%%%
\subsection{Batch Gradient Descent (BGD)}

Batch gradient descent (BGD), analyzed in Section~\ref{sec-grad-desc-basic}, iterates
\eq{
	x_{k+1} = x_k - \tau_k \nabla f(x_k)
}
For a convex objective with an $L$-Lipschitz gradient and a minimizer, the condition $0<\tau_{\min}\leq\tau_k\leq\tau_{\max}<2/L$ ensures convergence. Strong convexity strengthens this to a linear rate.

For a finite-sum objective, the gradient is
\eql{\label{eq-full-grad}
	\nabla f(x) = \frac{1}{n} \sum_{i=1}^n \nabla f_i(x)
}
Its cost is typically $O(np)$ when each $\nabla f_i$ can be evaluated in time linear in $p$.

For an ERM summand of the form~\eqref{eq-stochastic-erm}, Taylor expansion gives
\begin{align*}
	f_i(x+\epsilon) &= \ell(\dotp{a_i}{x} + \dotp{a_i}{\epsilon},y_i)
	= \ell(\dotp{a_i}{x},y_i) + \ell'(\dotp{a_i}{x},y_i)\dotp{a_i}{\epsilon} + o(\norm{\epsilon}) \\
	&= f_i(x) + \dotp{\ell'(\dotp{a_i}{x},y_i)a_i}{\epsilon} + o(\norm{\epsilon}),
\end{align*}
Here $\ell'(s,y)$ denotes differentiation with respect to the score $s$, with the response $y$ fixed. Consequently,
\eql{\label{eq-grad-formula}
	\nabla f_i(x) =  \ell'( \dotp{a_i}{x},y_i ) a_i.
}
For the logistic loss, one has 
\eq{
	\loss'(s,y)=-y \frac{e^{-sy}}{ 1+e^{-sy} }.
}



\begin{figure}
\centering
\begin{tabular}{cc}
\includegraphics[width=.45\linewidth]{ml/sgd/error-bgd-1} &
\includegraphics[width=.45\linewidth]{ml/sgd/error-bgd-2} \\
$f(x_k)$ & $\log_{10}(f(x_k)-f(x^\star))$ 
\end{tabular}
\caption{\label{fig-bgd}
Objective-error evolution for batch gradient descent in logistic classification.
}
\end{figure}



%%%%%%%%%%%%%%%%%%%%%%%%%%%%%%%%%%%%%%%%%%%%%%%%%%%%%%%%
\subsection{Stochastic Gradient Descent (SGD)}

\wrapf{ml/sgd/unbiased-grad}{Unbiased gradient estimate}
For very large $n$, computing the full gradient $\nabla f$ as in~\eqref{eq-full-grad} is prohibitive.  
%
SGD replaces the full gradient by the gradient of one summand $f_{i}$, with $i$ sampled uniformly. This gives an unbiased estimate:
\eql{\label{eq-unbiased-grad}
	\EE_{\ip}{ \nabla f_{\ip}(x) } = \nabla f(x)
}
where $\ip$ is a random variable distributed uniformly in $\{1,\ldots,n\}$.

\wrapf{ml/sgd/sgd-schematic}{Schematic view of SGD iterates}
Starting from $x_0$, stochastic gradient descent (SGD) iterates
\eq{
	x_{k+1} = x_k - \tau_k \nabla f_{i(k)}(x_k)
}
where, for each iteration index $k$, $i(k)$
is drawn independently of previous draws, uniformly from $\{1,\ldots,n\}$.
%
The iterates $x_{k+1}$ are random. We therefore study convergence to a minimizer of $f$ in a probabilistic sense, such as mean square or probability, together with the corresponding rates.

A batch-gradient step costs $O(np)$,
whereas an SGD step costs $O(p)$. This makes SGD attractive
when $n$ is large and repeated passes through the data are costly.
Redundancy between observations can sometimes yield
accurate predictions after only $k \ll n$ updates, before all
observations have been used.

The step-size schedule $\tau_k$ must balance two requirements: it
usually needs to decrease to suppress stochastic noise,
but not so quickly that the method stops making progress
before reaching a minimizer.


A common schedule has asymptotic behavior $\tau_k \sim k^{-1}$
as $k\rightarrow +\infty$. One such schedule is
\eql{\label{eq-stepsize-sgd}
	\tau_k \eqdef \frac{\tau_0}{1 + k/k_0}
}
where $\tau_0>0$ sets the initial step and $k_0>0$ controls how long the steps remain approximately constant.

Figure~\ref{fig-sgd-traject} illustrates minimization of $f_1(x)+f_2(x)$ for $x \in \RR$, with $f_1(x)=(x-1)^2$ and $f_2(x)=(x+1)^2$. The distribution of $x_k$ concentrates around the minimizer $x^\star=0$. The initial iterate $x_0$ is uniform on $[-1/2,1/2]$.

\begin{figure}
\centering
\includegraphics[width=.49\linewidth]{ml/sgd/sgd-histo} 
\includegraphics[width=.49\linewidth]{ml/sgd/sgd-trajectory}
\caption{\label{fig-sgd-traject}
Repeated SGD runs for $k\mapsto x_k\in\RR$. Left: the distribution of iterates at each iteration. Right: individual trajectories.
}
\end{figure}

The following theorem gives a mean-square convergence rate. The second-moment assumption is imposed along the iterates, not uniformly on all of $\RR^p$: a globally strongly convex function cannot have a globally bounded gradient.

\begin{thm}\label{thm-conv-sgd}
Let $f$ be differentiable and $\mu$-strongly convex, with minimizer $x^\star$. Let $\mathcal F_k$ contain the initial point and all randomness before the $k$th update, so $x_k$ is $\mathcal F_k$-measurable. Assume the stochastic gradient $G_k$ satisfies
\eq{\EE(G_k\mid\mathcal F_k)=\nabla f(x_k),\qquad
\EE(\norm{G_k}^2\mid\mathcal F_k)\leq C^2.}
For $x_{k+1}=x_k-\tau_kG_k$, $\tau_k=1/(\mu(k+2))$, and deterministic $x_0$,
\eql{\label{eq-rate-sgd}
\EE\norm{x_k-x^\star}^2\leq\frac{R}{k+2},\qquad
R=\max\{2\norm{x_0-x^\star}^2,C^2/\mu^2\}.
}
If $f$ additionally has an $L$-Lipschitz gradient, then
$\EE(f(x_k)-f(x^\star))\leq LR/(2(k+2))$.
\end{thm}
\begin{proof}
Strong convexity and $\nabla f(x^\star)=0$ imply
\eql{\label{eq-sgd-proof-1}
\dotp{\nabla f(x_k)}{x_k-x^\star}\geq\mu\norm{x_k-x^\star}^2.
}
Expand the squared update and take a conditional expectation, then average over the history. With $e_k=\EE\norm{x_k-x^\star}^2$,
\eql{\label{eq-sgd-proof-2}
e_{k+1}\leq(1-2\mu\tau_k)e_k+\tau_k^2C^2.
}
The claim holds at $k=0$ by the definition of $R$. Set $m=k+2\geq2$. Since $1-2/m\geq0$, the induction hypothesis gives
\eq{e_{k+1}\leq\left(1-\frac2m\right)\frac Rm+\frac R{m^2}
=\frac{R(m-1)}{m^2}\leq\frac R{m+1}.}
The objective bound follows from $L$-smoothness at the minimizer.
\end{proof}

The preceding SGD bound, and the averaging scheme below, give sublinear rates even when $f$ is strongly convex. By contrast, batch gradient descent has a linear rate under the assumptions of Theorem~\ref{thm-gradsec-non-strong-conv}.

Figure~\ref{fig-sgd} compares the objective values $f(x_k)$.
The black dashed curve shows batch gradient descent, with the iteration
axis rescaled to reflect the $n$-fold cost of a batch update.


\begin{figure}
\centering
\begin{tabular}{cc}
\includegraphics[width=.45\linewidth]{ml/sgd/error-sgd-1} &
\includegraphics[width=.45\linewidth]{ml/sgd/error-sgd-2} \\
$f(x_k)$ & $\log_{10}(f(x_k)-f(x^\star))$ 
\end{tabular}
\caption{\label{fig-sgd}
Objective-error evolution for SGD in logistic classification; the dashed curve shows batch gradient descent.
}
\end{figure}




%%%%%%%%%%%%%%%%%%%%%%%%%%%%%%%%%%%%%%%%%%%%%%%%%%%%%%%%
\subsection{Stochastic Gradient Descent with Averaging (SGA)}
\label{sec-sga}

Diminishing steps can slow SGD as
$\tau_k$ approaches zero.
%
Averaging reduces fluctuations in the estimate. Run SGD on auxiliary
iterates $( \tilde x_k)_k$,
\eq{
	 \tilde x_{k+1}=\tilde x_k - \tau_k \nabla f_{i(k)}(\tilde x_k)
}
and return their Cesàro average:
\eq{
	x_k \eqdef \frac{1}{k} \sum_{\ell=1}^k \tilde x_{\ell}.
}
This gives stochastic gradient descent with averaging (SGA),
also called averaged SGD.

The average can be maintained without storing the full trajectory,
using the recursive update
\eq{
	x_{k+1}=\frac{1}{k+1}\tilde x_{k+1}+\frac{k}{k+1}x_k.
}


A typical step-size schedule is
\eq{
	\tau_k \eqdef \frac{\tau_0}{1 + \sqrt{k/k_0}}.
}
These steps decrease more slowly, at rate $k^{-1/2}$.


Averaging can reduce sensitivity to the tuning parameters
$(k_0,\tau_0)$, although their choice still affects performance.

% <https://arxiv.org/pdf/1303.6149.pdf 

Bach proves that for logistic classification, 
averaging can adapt to local curvature under suitable regularity and step-size assumptions. This is a problem-dependent result, rather than a universal improvement of every constant over SGD.



%%%%%%%%%%%%%%%%%%%%%%%%%%%%%%%%%%%%%%%%%%%%%%%%%%%%%%%%
\subsection{Stochastic Averaged Gradient Descent (SAG)}

% https://arxiv.org/pdf/1309.2388
For a finite dataset with $n$ samples of total size $n \times p$,
one can trade additional memory for reduced gradient noise
by storing past gradient evaluations. This yields
the stochastic averaged gradient (SAG) method.

Store the latest gradient of each summand in a table $(G^i)_{i=1}^n$, requiring $O(np)$ memory. Maintain its average $g$ as an approximation to the full gradient.

Initialize $G^i=0$ for all $i$ and $g=0$, so that $g=n^{-1}\sum_iG^i$ is the stored average. At each iteration, update the table and then $x$:
\eq{
	x_{k+1} = x_k - \tau g
	\qwhereq
	\choice{
	h\leftarrow\nabla f_{i(k)}(x_k), \\
	g\leftarrow g+\frac1n(h-G^{i(k)}),   \\
	G^{i(k)} \leftarrow h.
	}
}
Because most entries were evaluated at earlier iterates, the table average is generally biased conditional on the current history. Its convergence analysis therefore differs from the unbiased-SGD proof above. SAG uses a fixed step size $\tau$. As for batch gradient descent,
choose $\tau$ on the scale of $1/L$,
where $L$ bounds the Lipschitz constants of the individual gradients $\nabla f_i$; the admissible numerical constant depends on the convergence theorem.

For smooth convex finite sums with a minimizer and a suitable step size, SAG admits an $O(1/k)$ expected objective bound for averaged iterates.
Furthermore, in the presence of strong convexity (for instance after adding a positive quadratic penalty to logistic regression), it has a linear convergence rate,
i.e. 
 \eq{
	\EE( f(x_k) ) - f(x^\star) = O\pa{ \rho^k },
}
for some $0 < \rho < 1$. 

These improvements exploit the finite-sum structure:
SAG stores information about each of the $n$ summands. SGD and averaged SGD
also extend beyond finite $n$ to general
expectations~\eqref{eq-min-int}.

Figure~\ref{fig-compariso-sgd} shows a comparison of SGD, SGA and SAG.

\begin{figure}
\centering
\includegraphics[width=.5\linewidth]{ml/sgd/sg-comparison}
\caption{\label{fig-compariso-sgd}
Evolution of $\log_{10}(f(x_k)-f(x^\star))$ for SGD, SGA and SAG.
}
\end{figure}



